\chapter{Самостоятельные задания}
Пишите промежуточные шаги. В заданиях с битовыми словами показывайте ведущие нули до указанной ширины. Изменения одного слова в разных подпунктах выполняйте независимо, если прямо не указана последовательность.

\section{Блок 1. Переводы}
\subsection*{З1. Из десятичной системы}\label{task:z1}
Переведите $214_{10}$ в основания 2, 8 и 16. Покажите деление с остатком хотя бы для одного основания и проверьте ответы развёрнутой записью.
\hyperref[ans:z1]{Ответ и решение.}

\subsection*{З2. В десятичную систему}\label{task:z2}
Переведите $1011011_2$, $273_8$ и $7D_{16}$ в десятичную систему. Укажите веса ненулевых разрядов.
\hyperref[ans:z2]{Ответ и решение.}

\subsection*{З3. Группировка}\label{task:z3}
Переведите $1110101101_2$ в основания 4, 8 и 16 без промежуточного перевода в десятичную систему. Покажите границы групп.
\hyperref[ans:z3]{Ответ и решение.}

\subsection*{З4. Обратный перевод и допустимость цифр}\label{task:z4}
\begin{enumerate}
    \item Переведите $3C7_{16}$ в двоичную и восьмеричную системы через группы битов.
    \item Какие записи допустимы: $102_2$, $19_8$, $2G_{16}$, $203_4$? Для каждой недопустимой записи назовите причину.
\end{enumerate}
\hyperref[ans:z4]{Ответ и решение.}

\section{Блок 2. Арифметика}
\subsection*{З5. Двоичное сложение и вычитание}\label{task:z5}
Выполните $11011_2+10110_2$ и $101101_2-1110_2$. Покажите переносы и заёмы; разрядность не ограничена.
\hyperref[ans:z5]{Ответ и решение.}

\subsection*{З6. Шестнадцатеричная арифметика}\label{task:z6}
Вычислите $2B7_{16}+8D_{16}$ и $300_{16}-B6_{16}$. Объясните заём через нулевой разряд во втором действии.
\hyperref[ans:z6]{Ответ и решение.}

\subsection*{З7. Умножение}\label{task:z7}
Вычислите $34_8\cdot6_8$ и $1011_2\cdot101_2$ в указанных основаниях. Проверьте в десятичной системе.
\hyperref[ans:z7]{Ответ и решение.}

\subsection*{З8. Частное и остаток}\label{task:z8}
Найдите частное и остаток от $325_8:6_8$ и $2D9_{16}:B_{16}$. Для каждой пары проверьте равенство $A=BQ+R$ и неравенство $0\le R<B$.
\hyperref[ans:z8]{Ответ и решение.}

\section{Блок 3. Диапазоны}
\subsection*{З9. Коды и значения}\label{task:z9}
Для 6-битных unsigned, прямого, обратного и дополнительного кодов укажите диапазон, количество битовых комбинаций и количество различных числовых значений. Объясните, почему последние два количества могут различаться.
\hyperref[ans:z9]{Ответ и решение.}

\subsection*{З10. Девять битов}\label{task:z10}
Для 9-битных unsigned и дополнительного кода укажите границы, количество значений и битовые записи минимального и максимального значений.
\hyperref[ans:z10]{Ответ и решение.}

\subsection*{З11. Минимальная разрядность}\label{task:z11}
Найдите минимальное число битов для каждого случая:
\begin{enumerate}
    \item unsigned со всеми значениями от 0 до 1000;
    \item дополнительный код со всеми значениями от $-300$ до 300;
    \item дополнительный код со всеми значениями от $-128$ до 128.
\end{enumerate}
Докажите, что на один бит меньше недостаточно.
\hyperref[ans:z11]{Ответ и решение.}

\subsection*{З12. Пограничные числа}\label{task:z12}
Какие из чисел $-128$, 127, 128 можно представить в каждом из 8-битных форматов: unsigned, прямой код, обратный код, дополнительный код, код со смещением 128? Для представимых чисел приведите hex-код, для остальных объясните ограничение.
\hyperref[ans:z12]{Ответ и решение.}

\section{Блок 4. Кодирование}
\subsection*{З13. Два знака, четыре кода}\label{task:z13}
Запишите $+46$ и $-46$ в 8-битных прямом, обратном, дополнительном кодах и коде со смещением 128. Дайте двоичные и шестнадцатеричные записи.
\hyperref[ans:z13]{Ответ и решение.}

\subsection*{З14. Одна запись --- пять интерпретаций}\label{task:z14}
Прочитайте 8-битное слово $D6_{16}$ как unsigned, прямой, обратный, дополнительный код и код со смещением 128.
\hyperref[ans:z14]{Ответ и решение.}

\subsection*{З15. Сложение смещённых кодов}\label{task:z15}
Числа $A=-40$, $B=17$ представлены в 8 битах со смещением 128. Найдите их коды, код суммы в том же формате и код суммы в дополнительном коде. Объясните, почему обычного сложения кодов недостаточно.
\hyperref[ans:z15]{Ответ и решение.}

\subsection*{З16. Особые значения дополнительного кода}\label{task:z16}
Запишите $-1$, 0 и $-128$ в 8-битном дополнительном коде. Примените к коду $-128$ операцию «инверсия плюс один» в 8 битах. Какой код получится? Можно ли считать его представлением числа $+128$ в этом формате?
\hyperref[ans:z16]{Ответ и решение.}

\section{Блок 5. Арифметика и флаги}
Во всех заданиях блока ширина 8 битов. Укажите код результата, его беззнаковое и знаковое значения, затем CF, OF, ZF, SF. Десятичные отрицательные операнды кодируйте дополнительным кодом. Для вычитания определяйте флаги SUB.

\subsection*{З17. Сложение разных знаков}\label{task:z17}
Выполните $52+(-19)$. Объясните, почему перенос и знаковое переполнение дают разные ответы.
\hyperref[ans:z17]{Ответ и решение.}

\subsection*{З18. Сложение одинаковых знаков}\label{task:z18}
Выполните $95+48$ и $(-100)+(-40)$. В обоих случаях сравните точную математическую сумму со знаковой интерпретацией сохранённого слова.
\hyperref[ans:z18]{Ответ и решение.}

\subsection*{З19. Две границы}\label{task:z19}
Сложите коды $F0_{16}+10_{16}$ и $7F_{16}+01_{16}$. Покажите, что ненулевой CF не требует ненулевого OF и наоборот.
\hyperref[ans:z19]{Ответ и решение.}

\subsection*{З20. Флаги вычитания}\label{task:z20}
Выполните $35-60$ и $(-110)-30$. Найдите младшие биты через сложение с дополнительным кодом вычитаемого, но CF определите для исходного вычитания.
\hyperref[ans:z20]{Ответ и решение.}

\section{Блок 6. Побитовые операции и маски}
\subsection*{З21. Четыре операции}\label{task:z21}
Пусть $x=AC_{16}$, $y=69_{16}$ --- 8-битные слова. Найдите AND, OR, XOR этих слов и NOT $x$. Ответы запишите в двоичной и шестнадцатеричной системах.
\hyperref[ans:z21]{Ответ и решение.}

\subsection*{З22. Отдельные биты}\label{task:z22}
Для исходного 8-битного $x=53_{16}$ независимо:
\begin{enumerate}
    \item установите бит 7;
    \item сбросьте бит 4;
    \item инвертируйте бит 0;
    \item проверьте, установлен ли бит 6;
    \item выделите младшие 3 бита.
\end{enumerate}
Для каждого действия укажите операцию, маску и результат.
\hyperref[ans:z22]{Ответ и решение.}

\subsection*{З23. Поля битов}\label{task:z23}
Для 8-битного $x=DA_{16}$:
\begin{enumerate}
    \item извлеките биты с номерами 4--6 как unsigned-число от 0 до 7;
    \item независимо сбросьте бит 3.
\end{enumerate}
Затем соберите отдельный байт, у которого старшая тетрада равна $A_{16}$, а младшая --- $5_{16}$. Покажите сдвиги и маски.
\hyperref[ans:z23]{Ответ и решение.}

\subsection*{З24. Замена поля}\label{task:z24}
У 8-битного $x=B6_{16}$ замените младшую тетраду на $9_{16}$, сохранив старшую. Запишите последовательность AND и OR. Объясните, почему одного OR с $09_{16}$ недостаточно.
\hyperref[ans:z24]{Ответ и решение.}

\section{Блок 7. Сдвиги и ширина слова}
\subsection*{З25. Логические сдвиги}\label{task:z25}
Для исходного 8-битного $x=6D_{16}$ независимо выполните логические сдвиги влево на 1 и на 3, вправо на 2. Запишите биты и hex-код. Для левых сдвигов сравните результат с точным беззнаковым произведением.
\hyperref[ans:z25]{Ответ и решение.}

\subsection*{З26. Арифметический сдвиг и деление}\label{task:z26}
Дано 8-битное $x=93_{16}$ в дополнительном коде. Найдите его значение. Независимо выполните арифметические сдвиги вправо на 1 и на 3, логический вправо на 3. Сравните арифметический сдвиг на 3 с делением исходного знакового числа на 8 в C.
\hyperref[ans:z26]{Ответ и решение.}

\subsection*{З27. Циклические сдвиги}\label{task:z27}
Для 8-битного $x=81_{16}$ независимо выполните циклический сдвиг влево на 1 и вправо на 2. Сравните первый результат с логическим сдвигом влево на 1.
\hyperref[ans:z27]{Ответ и решение.}

\subsection*{З28. Расширение и сужение}\label{task:z28}
\begin{enumerate}
    \item Расширьте 8-битное $D3_{16}$ до 16 битов: сначала как unsigned, затем как дополнительный код. Укажите значения.
    \item Сузьте 16-битное unsigned-число $0137_{16}$ до 8 битов, отбросив старшие биты. Сохранится ли его значение?
\end{enumerate}
\hyperref[ans:z28]{Ответ и решение.}

\section{Блок 8. Комбинированные задачи}
\subsection*{З29. От кода к вычислению}\label{task:z29}
Прочитайте 8-битное $E9_{16}$ в дополнительном коде и выполните его знаковое расширение до 16 битов. Затем к \textbf{исходному 8-битному слову} прибавьте 8 в 8-битной сетке. Укажите код результата, знаковое значение и четыре флага.
\hyperref[ans:z29]{Ответ и решение.}

\subsection*{З30. Найдите ошибки в утверждениях}\label{task:z30}
Исправьте каждое утверждение и приведите объяснение или контрпример:
\begin{enumerate}
    \item «Слово \texttt{1111} всегда означает $-1$».
    \item «Если при сложении CF=1, то обязательно OF=1».
    \item «В 8 битах инверсия $0F_{16}$ равна $00_{16}$, потому что инвертируются только четыре записанные единицы».
\end{enumerate}
\hyperref[ans:z30]{Ответ и решение.}

\subsection*{З31. Найдите ошибку в округлении}\label{task:z31}
Студент записал: «$-3$ в 8 битах имеет код $FD_{16}$. Арифметический сдвиг вправо на 1 равен делению на 2 в C, поэтому результат $-1$, код $FF_{16}$». Найдите первый неверный шаг и покажите правильные биты сдвига.
\hyperref[ans:z31]{Ответ и решение.}

\subsection*{З32. Другая ширина и другое смещение}\label{task:z32}
В 7-битном коде со смещением 64 представлены $A=-25$ и $B=38$. Найдите диапазон формата, коды операндов и код суммы. Что получится, если просто сложить коды, оставить 7 младших битов и прочитать их в исходном формате? Запишите правильную сумму также в 7-битном дополнительном коде.
\hyperref[ans:z32]{Ответ и решение.}
